\section{Experimental Setup}
\label{sec:experimental-setup}

\paragraph{Models and benchmark.}
We evaluate Security Cuff in the deployment setting of
OpenVLA-7B~\citep{openvla} on the LIBERO benchmark~\citep{libero}.
Our target evaluation scope covers the full LIBERO suite, including
\texttt{LIBERO-Spatial}, \texttt{LIBERO-Object},
\texttt{LIBERO-Goal}, and \texttt{LIBERO-Long}.
All experiments are organized at the rollout level, because the system
must decide whether a trajectory should be continued, escalated, or
intervened upon during execution rather than merely classify a single
frame or trigger instance.

\paragraph{Rollout pool and attack families.}
The evaluation pool contains both benign and adversarial rollouts.
Benign rollouts are divided into \texttt{task\_success} and
\texttt{task\_fail}.
Adversarial rollouts are grouped into three attack families:
\texttt{attack\_GoBA}, \texttt{attack\_Drop}, and
\texttt{attack\_State}.
GoBA-style attacks provide physically triggered backdoor behavior via
object- or scene-dependent triggers~\citep{goba}.
DropVLA-family attacks provide text, visual, and joint-trigger
conditions generated from the official data release and fine-tuning
pipeline.
State-based attacks provide a complementary family in which the trigger
is tied to rollout state rather than an explicit visual artifact
observable in the scene~\citep{statebackdoor}.
Unless otherwise noted, our primary task is the pooled
\texttt{success\_vs\_other} protocol, which merges task failures and
all attack families into a single risky class.

% Pending execution note:
% The intended final protocol includes GoBA, DropVLA, and StateBackdoor
% under a unified pooled rollout label space. If some attack families
% are still incomplete in the current draft, keep this paragraph as the
% target protocol and only annotate temporary gaps in comments.

\paragraph{Baselines.}
We compare the full two-layer system against both system-level and
fast-layer baselines.
At the system level, we compare against simplified variants that remove
either routing or verification, so that the effect of the dual-layer
architecture can be isolated from the effect of any single detector.
At the fast-layer level, we compare against SAFE-style failure
detectors~\citep{safe}, including MLP and LSTM probes, as well as
distance-based baselines such as cosine $k$-NN and Mahalanobis
detectors.
These baselines test three distinct questions:
whether rollout risk can be detected from hidden features at all,
whether sequential modeling helps relative to static probing, and
whether the system gain comes from the dual-layer architecture rather
than from a single stronger feature-space classifier.

\paragraph{Metrics.}
We report both detection quality and deployment cost.
For trajectory-level detection, we report AUROC, balanced accuracy, and
recall at a target false positive rate.
For early warning, we report the mean first detection step.
For system-level defense, we report clean success rate and attack
success rate after defense, since a runtime safety system must preserve
normal behavior while suppressing harmful rollouts.
For deployment cost, we report latency, memory overhead, and routing
rate to the slow layer.
Together, these metrics evaluate effectiveness, timeliness, and
practicality under the formulation in
Section~\ref{sec:problem-formulation}.

\paragraph{Main tables and ablations.}
Our main evidence block is a pair of system-level tables in the main
paper: a whole-system comparison table and a system overhead table.
The whole-system table reports defense effectiveness, clean success
preservation, and attack suppression across LIBERO suites.
The overhead table reports latency and memory cost of the deployed
defense stack.
Fast-layer ablations are reported separately to analyze raw versus
residual features, SAFE-style baselines, and the effect of pooled
training versus transfer from failure-only supervision.
Additional cross-attack experiments evaluate mixed-family robustness
and leave-one-family-out generalization.
